\section{Typing Rules}

The base typing rules for programs, statements, and expressions are given in
\Cref{fig:tr-prog,fig:tr-stmt,fig:tr-expr}.
They use the following notation:

\vspace{1em}
\begin{tabular}{r | l}
    $\Gamma,\Phi \vdash e : \tau$
      & expression $e$ in context $\Gamma,\Phi$ has type $\tau$ \\
    $\Gamma,\Phi \vdash s\ \triangleright\ \Gamma_s,\Phi_s$
      & statement $s$ in context $\Gamma,\Phi$ produces context $\Gamma_s,\Phi_s$ \\
    $\Gamma,\Phi \vdash p \blacktriangleright \tau$
      & program $p$ in context $\Gamma,\Phi$ outputs a value of type $\tau$ \\
    $\Gamma[x \mapsto \tau]$
      & variable context $\Gamma$ extended with $x : \tau$ \\
    $\Phi[f \mapsto (\sigma \to \tau)]$
      & function context $\Phi$ extended with $f : \sigma \to \tau$ \\
    $(x : \tau) \in \Gamma$
      & $x$ is bound to type $\tau$ in $\Gamma$ \\
    $(f : \sigma \to \tau) \in \Phi$
      & $f$ is bound in $\Phi$ \\
\end{tabular}

\begin{figure}
    \small

    % Program --------------------------------------------

    \begin{prooftree}
        \AxiomC{$\Gamma,\Phi \vdash e : \tau$}

        \RightLabel{[\textsc{P-Ret}]}

        \UnaryInfC{$\Gamma,\Phi \vdash e \blacktriangleright \tau$}
    \end{prooftree}

    \begin{prooftree}
        \AxiomC{$\Gamma,\Phi \vdash s\ \triangleright\ \Gamma_s,\Phi_s$}
        \AxiomC{$\Gamma_s,\Phi_s \vdash p\ \blacktriangleright\ \tau$}

        \RightLabel{[\textsc{P-Seq}]}

        \BinaryInfC{$\Gamma,\Phi \quad\vdash\quad s \texttt{;}\ p \quad\blacktriangleright\quad \tau$}
    \end{prooftree}

\caption{Typing rules for programs}
\label{fig:tr-prog}
\end{figure}

\begin{figure}
    \small
    % Statements -----------------------------------------

    % SLet. Stmt ::= "val" Ident "=" Exp ;
    \begin{prooftree}
        \AxiomC{$\Gamma,\Phi \vdash e : \tau$}

        \RightLabel{[\textsc{S-Val}]}

        \UnaryInfC{$\Gamma,\Phi \quad\vdash\quad \texttt{val}\ x\ \texttt{=}\ e \quad\triangleright\quad \Gamma[x \mapsto \tau],\Phi$}
    \end{prooftree}

    % SFun. Stmt ::= "fun" Ident "(" Ident ":" Type ") =" Exp ;
    \begin{prooftree}
        \AxiomC{$\Gamma[x \mapsto \sigma],\Phi \vdash e : \tau$}

        \RightLabel{[\textsc{S-Fun}]}

        \UnaryInfC{$\Gamma,\Phi \quad\vdash\quad \texttt{fun}\ f\ \texttt{(}\ x\ \texttt{:}\ \sigma\ \texttt{) =}\ e \quad\triangleright\quad \Gamma,\Phi[f \mapsto (\sigma \to \tau)]$}
    \end{prooftree}

\caption{Typing rules for statements}
\label{fig:tr-stmt}
\end{figure}

\begin{figure}
    \small

    % Primitives -----------------------------------------

    \begin{multicols}{2}
        % EInt. Exp7  ::= Integer;
        \begin{prooftree}
            \AxiomC{$i \in \mathbb{Z}$}

            \RightLabel{[\textsc{E-Int}]}

            \UnaryInfC{$\vdash i : int$}
        \end{prooftree}

        \columnbreak

        % ETrue.  Exp3  ::= "True";
        % EFalse. Exp3  ::= "False";
        \begin{prooftree}
            \AxiomC{$b \in \{\texttt{True}, \texttt{False}\}$}

            \RightLabel{[\textsc{E-Bool}]}

            \UnaryInfC{$\vdash b : bool$}
        \end{prooftree}
        
    \end{multicols}

    % Arithmetic -----------------------------------------

    % EMul. Exp6  ::= Exp6 "*" Exp7;
    % EDiv. Exp6  ::= Exp6 "/" Exp7;
    % EAdd. Exp5  ::= Exp5 "+" Exp6;
    % ESub. Exp5  ::= Exp5 "-" Exp6;
    \begin{prooftree}
        \AxiomC{$\Gamma,\Phi \vdash e_1 : int$}
        \AxiomC{$\Gamma,\Phi \vdash e_2 : int$}
        \AxiomC{$\star \in \{\texttt{*}, \texttt{/}, \texttt{+}, \texttt{-}\}$}

        \RightLabel{[\textsc{E-Arith}]}

        \TrinaryInfC{$\Gamma,\Phi \vdash e_1 \star e_2 : int$}
    \end{prooftree}

    % Logic ----------------------------------------------

    % ENot.   Exp3  ::= "!" Exp3;
    \begin{prooftree}
        \AxiomC{$\Gamma,\Phi \vdash e : bool$}

        \RightLabel{[\textsc{E-Not}]}

        \UnaryInfC{$\Gamma,\Phi \vdash \texttt{!}e : bool$}
    \end{prooftree}

    % EAnd.   Exp2  ::= Exp2 "&&" Exp3;
    % EOr.    Exp1  ::= Exp1 "||" Exp2;
    \begin{prooftree}
        \AxiomC{$\Gamma,\Phi \vdash e_1 : bool$}
        \AxiomC{$\Gamma,\Phi \vdash e_2 : bool$}
        \AxiomC{$\star \in \{\texttt{\&\&}, \texttt{||}\}$}

        \RightLabel{[\textsc{E-Logic}]}

        \TrinaryInfC{$\Gamma,\Phi \vdash e_1 \star e_2 : bool$}
    \end{prooftree}

    % Comparisons ----------------------------------------

    % EEq.  Exp  ::= Exp "==" Exp1;
    \begin{prooftree}
        \AxiomC{$\Gamma,\Phi \vdash e_1 : \tau$}
        \AxiomC{$\Gamma,\Phi \vdash e_2 : \tau$}
        \AxiomC{$\tau \in \{int, bool\}$}

        \RightLabel{[\textsc{E-Eq}]}

        \TrinaryInfC{$\Gamma,\Phi \vdash e_1\ \texttt{==}\ e_2 : bool$}
    \end{prooftree}

    % ELt.  Exp4 ::= Exp4 "<" Exp5;
    % EGt.  Exp4 ::= Exp4 ">" Exp5;
    % ELeq. Exp4 ::= Exp4 "<=" Exp5;
    % EGeq. Exp4 ::= Exp4 ">=" Exp5;
    \begin{prooftree}
        \AxiomC{$\Gamma,\Phi \vdash e_1 : int$}
        \AxiomC{$\Gamma,\Phi \vdash e_2 : int$}
        \AxiomC{$\star \in \{\texttt{<}, \texttt{>}, \texttt{<=}, \texttt{>=}\}$}

        \RightLabel{[\textsc{E-Comp}]}

        \TrinaryInfC{$\Gamma,\Phi \vdash e_1 \star e_2 : bool$}
    \end{prooftree}

    % Control flow ---------------------------------------

    % EIf.  Exp8 ::= "if" Exp "then" Exp "else" Exp;
    \begin{prooftree}
        \AxiomC{$\Gamma,\Phi \vdash cond : bool$}
        \AxiomC{$\Gamma,\Phi \vdash e_1 : \tau$}
        \AxiomC{$\Gamma,\Phi \vdash e_2 : \tau$}

        \RightLabel{[\textsc{E-If}]}

        \TrinaryInfC{$\Gamma,\Phi \vdash \texttt{if}\ cond\ \texttt{then}\ e_1\ \texttt{else}\ e_2 : \tau$}
    \end{prooftree}

    % Let bindings ---------------------------------------

    % ELet. Exp8 ::= "let" Ident "=" Exp "in" Exp;
    \begin{prooftree}
        \AxiomC{$\Gamma,\Phi \vdash e : \sigma$}
        \AxiomC{$\Gamma[x \mapsto \sigma],\Phi \vdash body : \tau$}

        \RightLabel{[\textsc{E-Let}]}

        \BinaryInfC{$\Gamma,\Phi \quad\vdash\quad \texttt{let}\ x\ \texttt{=}\ e\ \texttt{in}\ body \quad:\quad \tau$}
    \end{prooftree}

    % EVar. Exp8 ::= Ident;
    \begin{prooftree}
        \AxiomC{$(x : \tau) \in \Gamma$}

        \RightLabel{[\textsc{E-Var}]}

        \UnaryInfC{$\Gamma,\Phi \vdash x : \tau$}
    \end{prooftree}

    % Functions ------------------------------------------

    % EApp. Exp8 ::= Exp8 Exp7;
    \begin{prooftree}
        \AxiomC{$(f : \sigma \to \tau) \in \Phi$}
        \AxiomC{$\Gamma,\Phi \vdash e : \sigma$}

        \RightLabel{[\textsc{E-App}]}

        \BinaryInfC{$\Gamma,\Phi \vdash f\ e : \tau$}
    \end{prooftree}
        
    \caption{Typing rules for expressions}
    \label{fig:tr-expr}
\end{figure}


% ---------------------------------------------------------------------------
% Extended rules for the ownership type system
% ---------------------------------------------------------------------------

\subsection*{Extended Context}

The variable context $\Gamma$ is extended to track ownership and borrow state.
Each entry maps a variable $x$ to a tuple $(\tau,\,\mu,\,\omega,\,b,\,m,\,\rho)$:

\vspace{0.5em}
\begin{tabular}{c | l}
    $\tau$   & type \\
    $\mu \in \{\mathtt{imm},\mathtt{mut}\}$  & mutability \\
    $\omega \in \{\checkmark, \times\}$ & owned ($\checkmark$) or moved ($\times$) \\
    $b \in \mathbb{N}$ & active immutable borrow count \\
    $m \in \{0,1\}$   & active mutable borrow count (exclusive) \\
    $\rho \in \{\varepsilon, x\}$ & borrow source ($\varepsilon$ if not a borrow) \\
\end{tabular}
\vspace{0.5em}

\noindent
We write $b_x$ and $m_x$ for the borrow counts of $x$ in the current $\Gamma$.
The predicate $\mathtt{Copy}(\tau)$ holds for $\tau \in \{\mathtt{int},\mathtt{bool},\,\&T\}$.
The modified expression judgment $\Gamma,\Phi \vdash e : \tau \mid \Gamma'$
indicates that evaluating $e$ may update $\Gamma$ (e.g.\ by marking a moved
variable); for \texttt{Copy} types $\Gamma' = \Gamma$ always.

The rules in \Cref{fig:tr-prog,fig:tr-stmt,fig:tr-expr} give the base language.
The following rules add ownership, borrowing, NLL, and thread-safe spawn.
Where a base rule is superseded (E-Var, S-Val, P-Seq), the new rule replaces it.

% ---------------------------------------------------------------------------
\subsection*{Phase 1 --- Affine Variable Access}
% ---------------------------------------------------------------------------

\textsc{E-Var} is replaced by two rules depending on whether the type is
\texttt{Copy}.
A \texttt{Copy} variable can be read any number of times without consuming it;
a non-\texttt{Copy} variable is \emph{moved} on first read: ownership is
transferred and the variable is marked as consumed.

\vspace{0.5em}
\begin{multicols}{2}

    \begin{prooftree}
        \AxiomC{$(x : \tau,\,\checkmark) \in \Gamma$}
        \AxiomC{$\mathtt{Copy}(\tau)$}
        \RightLabel{[\textsc{E-Var-Copy}]}
        \BinaryInfC{$\Gamma,\Phi \vdash x : \tau \mid \Gamma$}
    \end{prooftree}

    \columnbreak

    \begin{prooftree}
        \AxiomC{$(x : \tau,\,\checkmark) \in \Gamma$}
        \AxiomC{$\neg\,\mathtt{Copy}(\tau)$}
        \AxiomC{$b_x = m_x = 0$}
        \RightLabel{[\textsc{E-Var-Move}]}
        \TrinaryInfC{$\Gamma,\Phi \vdash x : \tau \mid \Gamma[x.\omega \leftarrow \times]$}
    \end{prooftree}

\end{multicols}

The two statement rules for let-bindings replace \textsc{S-Val} and carry
through the updated context $\Gamma'$ produced by evaluating the expression
(which may have consumed a non-\texttt{Copy} variable).
\textsc{S-LetMut} additionally marks the new binding as mutable.

\vspace{0.5em}
\begin{multicols}{2}

    \begin{prooftree}
        \AxiomC{$\Gamma,\Phi \vdash e : \tau \mid \Gamma'$}
        \RightLabel{[\textsc{S-LetImm}]}
        \UnaryInfC{$\Gamma,\Phi \;\vdash\; \mathtt{let}\ x = e
          \;\triangleright\; \Gamma' \cup \{x : (\tau,\,\mathtt{imm},\,\checkmark)\},\,\Phi$}
    \end{prooftree}

    \columnbreak

    \begin{prooftree}
        \AxiomC{$\Gamma,\Phi \vdash e : \tau \mid \Gamma'$}
        \RightLabel{[\textsc{S-LetMut}]}
        \UnaryInfC{$\Gamma,\Phi \;\vdash\; \mathtt{let\ mut}\ x = e
          \;\triangleright\; \Gamma' \cup \{x : (\tau,\,\mathtt{mut},\,\checkmark)\},\,\Phi$}
    \end{prooftree}

\end{multicols}

\vspace{0.5em}
Reassignment restores ownership: the old borrow-of link is cleared and
$x.\omega$ returns to $\checkmark$.
Assignment is only permitted when there are no active borrows of $x$.

\begin{prooftree}
    \AxiomC{$(x : \tau,\,\mathtt{mut},\,\checkmark,\,b_x{=}0,\,m_x{=}0) \in \Gamma$}
    \AxiomC{$\Gamma,\Phi \vdash e : \tau \mid \Gamma'$}
    \RightLabel{[\textsc{S-Assign}]}
    \BinaryInfC{$\Gamma,\Phi \quad\vdash\quad x = e
      \quad\triangleright\quad \Gamma'[x.\omega \leftarrow \checkmark,\; x.\rho \leftarrow \varepsilon],\,\Phi$}
\end{prooftree}

% ---------------------------------------------------------------------------
\subsection*{Phase 3A --- Immutable Borrows}
% ---------------------------------------------------------------------------

A borrow expression $\&x$ returns the address of $x$ without consuming it.
The corresponding let-binding increments $x$'s immutable borrow counter and
records that $r$ borrows from $x$.
A mutable borrow of $x$ must not already be active.

\vspace{0.5em}
\begin{multicols}{2}

    \begin{prooftree}
        \AxiomC{$(x : \tau,\,\checkmark,\,m_x{=}0) \in \Gamma$}
        \RightLabel{[\textsc{E-Ref}]}
        \UnaryInfC{$\Gamma,\Phi \vdash \&x : \&\tau \mid \Gamma$}
    \end{prooftree}

    \columnbreak

    \begin{prooftree}
        \AxiomC{$(x : \tau,\,\checkmark,\,m_x{=}0) \in \Gamma$}
        \RightLabel{[\textsc{S-LetBorrow}]}
        \UnaryInfC{$\Gamma,\Phi \;\vdash\; \mathtt{let}\ r = \&x
          \;\triangleright\; \Gamma[b_x \mathrel{+}= 1] \cup \{r : (\&\tau,\,\checkmark,\,\rho{=}x)\},\,\Phi$}
    \end{prooftree}

\end{multicols}

\vspace{0.5em}
Dereferencing an immutable reference returns the type of the referent.
The reference itself is \texttt{Copy} so it is not consumed.

\begin{prooftree}
    \AxiomC{$(r : \&\tau,\,\checkmark) \in \Gamma$}
    \RightLabel{[\textsc{E-Deref-Imm}]}
    \UnaryInfC{$\Gamma,\Phi \vdash *r : \tau \mid \Gamma$}
\end{prooftree}

% ---------------------------------------------------------------------------
\subsection*{Phase 3B --- Mutable Borrows}
% ---------------------------------------------------------------------------

A mutable borrow requires the variable to be declared mutable and to have
no existing borrows of either kind.
Only one mutable borrow may be active at a time.
Writing through a mutable reference ($*r = e$) checks that $e$ has the
referent type.

\vspace{0.5em}
\begin{multicols}{2}

    \begin{prooftree}
        \AxiomC{$(x : \tau,\,\mathtt{mut},\,\checkmark,\,b_x{=}m_x{=}0) \in \Gamma$}
        \RightLabel{[\textsc{S-LetMutBorrow}]}
        \UnaryInfC{$\Gamma,\Phi \;\vdash\; \mathtt{let}\ r = \&\!\mathtt{mut}\ x
          \;\triangleright\; \Gamma[m_x \mathrel{+}= 1] \cup \{r : (\&\!\mathtt{mut}\ \tau,\,\checkmark,\,\rho{=}x)\},\,\Phi$}
    \end{prooftree}

    \columnbreak

    \begin{prooftree}
        \AxiomC{$(r : \&\!\mathtt{mut}\ \tau,\,\checkmark) \in \Gamma$}
        \AxiomC{$\Gamma,\Phi \vdash e : \tau \mid \Gamma'$}
        \RightLabel{[\textsc{S-DerefAssign}]}
        \BinaryInfC{$\Gamma,\Phi \;\vdash\; *r = e \;\triangleright\; \Gamma',\,\Phi$}
    \end{prooftree}

\end{multicols}

% ---------------------------------------------------------------------------
\subsection*{Phase 3C --- Non-Lexical Lifetimes}
% ---------------------------------------------------------------------------

The base \textsc{P-Seq} rule is replaced by \textsc{P-Seq-NLL}, which releases
borrows that are no longer live before checking the remainder of the program.
The function $\mathrm{live}(\bar{p})$ returns the set of identifiers mentioned
syntactically in the remaining program $\bar{p}$.
The function $\mathrm{release}(\Gamma, L)$ decrements the borrow counter of
any variable whose borrow-holder is absent from $L$ and clears the
borrow-holder's $\rho$ field.

\begin{prooftree}
    \AxiomC{$\Gamma,\Phi \vdash s \;\triangleright\; \Gamma',\Phi'$}
    \AxiomC{$\Gamma'' = \mathrm{release}(\Gamma',\, \mathrm{live}(\bar{p}))$}
    \AxiomC{$\Gamma'',\Phi' \vdash \bar{p} \;\blacktriangleright\; \tau$}
    \RightLabel{[\textsc{P-Seq-NLL}]}
    \TrinaryInfC{$\Gamma,\Phi \quad\vdash\quad s\,;\,\bar{p} \quad\blacktriangleright\quad \tau$}
\end{prooftree}

% ---------------------------------------------------------------------------
\subsection*{Phase 4A --- Explicit Lifetime Annotations}
% ---------------------------------------------------------------------------

A lifetime-generic function declares lifetime parameters $\bar{\ell}$ and may
return a reference whose lifetime must be named among $\bar{\ell}$.
The call site rule \textsc{E-Call-Lt} resolves which argument provides the
return borrow by matching the return lifetime to a parameter's annotation.

\begin{prooftree}
    \AxiomC{$\ell \in \bar{\ell}\ \text{if}\ \mathit{retTy} \in \{\&^\ell\tau,\,\&\!\mathtt{mut}^\ell\!\tau\}$}
    \AxiomC{$\Gamma',\Phi \vdash \mathit{body} : \mathit{retTy} \mid \Gamma''$}
    \RightLabel{[\textsc{S-FunLt}]}
    \BinaryInfC{$\Gamma,\Phi \;\vdash\;
      \mathtt{fn}\ f\langle\bar{\ell}\rangle(\bar{p}) \to \mathit{retTy}\ \{\,\mathit{body}\,\}
      \;\triangleright\; \Gamma,\,\Phi[f \mapsto \langle\bar{\ell}\rangle(\bar{p}) \to \mathit{retTy}]$}
\end{prooftree}

\noindent
where $\Gamma' = \Gamma$ extended with bindings for parameters $\bar{p}$.

% ---------------------------------------------------------------------------
\subsection*{Phase 4B --- Thread-Safe Spawn}
% ---------------------------------------------------------------------------

A \texttt{spawn} block may only capture variables from the enclosing scope
whose type is \texttt{Copy}.
This statically prevents data races: \texttt{Copy} types have value semantics
(no aliasing), so sharing them between threads is always safe.

\begin{prooftree}
    \AxiomC{$\forall y \in \mathrm{fv}(\mathit{body}) \cap \mathrm{dom}(\Gamma):\;
             \mathtt{Copy}(\Gamma(y).\tau)$}
    \AxiomC{$\Gamma,\Phi \vdash \mathit{body}$}
    \RightLabel{[\textsc{S-Spawn}]}
    \BinaryInfC{$\Gamma,\Phi \quad\vdash\quad
      \mathtt{spawn}\ \{\,\mathit{body}\,\}
      \quad\triangleright\quad \Gamma,\,\Phi$}
\end{prooftree}

\noindent
Note that immutable references (\texttt{\&T}) are \texttt{Copy} and may be
shared; mutable references (\texttt{\&mut T}) are not \texttt{Copy} and are
therefore rejected at the spawn site.
